
\begin{definition}[Absolute value]
    For all $x \in \mathbb{R}$, we define its \textbf{absolute value} $|x|$ as 
    \[
        \left| x \right| = 
        \begin{cases}
            x, \text{ if } x \geq 0,\\
            -x, \text{ if } x < 0,
        \end{cases}
    \]
    
\end{definition}

\begin{proposition}
    Let $x, y \in \mathbb{R}, a \geq 0$, then 
    \begin{enumerate}
        \item $|x| \geq 0$
        \item $|x| \leq a \iff -a \leq x \leq a$
        \item $|xy| = |x| \cdot |y|$
        \item (\text{triangle inequality}) $|x + y| \leq |x| + |y|$
    \end{enumerate}
\end{proposition}

\begin{proposition}
    Let $x, y \in \mathbb{R}$ such that $x \leq y + \epsilon, \forall \epsilon > 0$, then $x \leq y$.
\end{proposition}

\begin{proof}
    We prove by contrapositive. Assume that $x < y$, WTS that $\exists \epsilon > 0, y + \epsilon < x$.
    \medskip

    Pick 
    \[
        \epsilon := \frac{x - y}{2}
    \]
    
    Then 
    \begin{align*}
        y + \epsilon &= y + \frac{x - y}{2} \\
        &= \frac{x+y}{2} \text{ by (D), (A2), (A3)} \\
        & < \frac{x + x}{2} \text{ by (O3)} \\
        & < x
    \end{align*}
\end{proof}

\subsection{Completeness Axiom / LUB property}

\begin{definition}[Bounded from above]
    Let $S \subset \mathbb{R}$, we say $S$ is \textbf{bounded from above} 
    if $\exists m \in \mathbb{R}$ s.t. $s \leq m, \forall s \in S$. 
    Such an $m$ is called an \textbf{upper bound} of $S$.
\end{definition}

\begin{definition}[Bounded]
    Let $S \subset \mathbb{R}$, we say $S$ is \textbf{bounded} 
    if it is bounded from both above and below.
\end{definition}


\begin{example}
    $[0, 1]:= \{ x \in \mathbb{R}, 0 \leq x \leq 1 \} $  has upper bounds 1, 1.5, 2, etc.
\end{example}

\begin{example}
    $\mathbb{N}$ is not bounded from above.
\end{example}

\begin{definition}[Maximum]
    If an upper bound $m$ of $S$ is an element of $S$, then $m$ is called the \textbf{maximum} of $S$, denoted 
    \[
        \max S
    \]
\end{definition}

\begin{example}
    $S = [ 0, 1 \} $ does not have a maximum, i.e. $\max S$ DNE.
\end{example}
\begin{proof}
    $\forall m \in S$, take $\epsilon = \frac{1}{2} (1 - m)$. Then $m + \epsilon < 1, m + \epsilon > m$. Hence $m$ is not an UB. Hence any $m$ cannot be UB.
\end{proof}

\begin{definition}[Supremum / least upper bound]
    Let $S \subset \mathbb{R}$ be bounded from above. Then $m \in \mathbb{R}$ is called the \textbf{supremum}, or \textbf{least upper bound} if 
    \begin{enumerate}
        \item \label{def:sup-1} $m$ is an upper bound of $S$ 
        \item \label{def:sup-2} $\forall m^{\prime} \in \mathbb{R}$, if $m^{\prime} < m$ then $m^{\prime}$ is not an upper bound of $S$.
    \end{enumerate} 
\end{definition}

\begin{remark}
    $\sup [0, 1\} = 1$. 
\end{remark}

\begin{proof}
    Suppose $\sup [0, 1) \neq 1$. 
    Since $1$ is an upper bound and condition \ref{def:sup-1} holds, 
    then condition \ref{def:sup-2} must fail. 
    Suppose $m < 1$ is an upper bound, 
    but $ \frac{1}{2}(m + 1) \in S, \frac{1}{2}(m +1) > m$.
\end{proof}

\begin{remarks}
    \hfill

    \begin{enumerate}
        \item $\sup S$ is unique. 
        \item $\sup S$ may not be an element of $S$. 
        \item If $\max S$ exists, then $\max S = \sum S$
    \end{enumerate}
\end{remarks}

\begin{example}
    $S = \emptyset$, then any $x \in \mathbb{R}$ is an upper bound and lower bound of $S$, but $\max S, \min S, \sup S, \inf S$ do not exist.
\end{example}

\begin{theorem}
    $\forall x, y \in \mathbb{R}, x < y$, $\exists r \in \mathbb{Q}$ s.t. 
    $x < r < y$.
\end{theorem}


\begin{theorem}[Completeness Axiom]
    If $S \subset \mathbb{R}, S \neq \emptyset$, and $S$ bounded from above, then $\sup S$ exists.
\end{theorem}

\begin{corollary}
    If $S \subset \mathbb{R}, S \neq \emptyset$, and $S$ bounded from below, then $\inf S$ exists.
\end{corollary}

\begin{proof}[Proof of CA $\iff$ corollary]
    Assume CA. 
    By CA, we know the existence of $\sup S$. 
        WTS that $\inf S$ exists, i.e. $\exists M \in \mathbb{R}$,
        $\forall s \in S, M \leq s$, and 
        $\forall M^{\prime} > M, \exists s \in S$ s.t. 
        $M \leq S < M^{\prime}$. 
        \medskip 

        Consider $S^{\prime} = \{ -x, x \in S \} $. 
        Since $S$ nonempty, $S^{\prime}$ is nonempty. 
        WTS $S^{\prime}$ bounded from above, i.e. $\exists M$ 
        s.t. $\forall s \in S^{\prime}$, $M \geq s$. 
        \medskip 

        Since $S$ bounded from below, $\exists M \in \mathbb{R}, \forall s \in S$,
        $s \geq M$. Pick arbitrary $s \in S^{\prime}, -s \in S \implies M \leq -s \iff -M \geq s$.
        Hence $S^{\prime}$  bounded from above. By CA, $\sup S^{\prime}$ exists, say $M_2$. 
        \medskip 

        We claim that $\inf S = -M_2$. 
        \begin{itemize}
            \item We first show $\forall s \in S, s \geq -M_2$. Fix $s \in S$ arbitrary, 
            \begin{align*}
                s \in S & \implies -s \in S^{\prime} \\
                & \implies -s \leq M_2 \\
                & \implies s \geq -M_2 \\
            \end{align*}
            \item We then show that any $M^{\prime} > -M_2$ is not a lower bound. 
            Take arbirary $M^{\prime} > -M_2$
            \begin{align*}
                M^{\prime} > -M_2 & \implies - M^{\prime} < M_2 \\
                & \implies \exists s \in S^{\prime} \text{ s.t } -M^{\prime} < s \leq M_2 \\
                & \implies M^{\prime} > -s \geq M_2
            \end{align*}
            Since $s \in S^{\prime} \iff -s \in S$, $M^{\prime}$ is not a lower bound of $S$.
        \end{itemize} 
\end{proof}

\begin{remark}
    The completeness axiom does not hold in $\mathbb{Q}$.
\end{remark}

\begin{proof}
    Take $S = \{ x \in \mathbb{Q}: x \geq 0, x^2 < 2 \} $. 

    We claim that any $m \in \mathbb{Q}$ that upper bounds 
    $S$ cannot be its supremum. 
    \remind{Exercise}.
\end{proof}

\begin{remark}[Implications of CA]
    \medskip
    \begin{enumerate}
        \item $\mathbb{Q}$ dense in $\mathbb{R}$.
        \item $\mathbb{R} \setminus \mathbb{Q}$ dense in $\mathbb{R}$.
        \item $x^2 = 2$ has solution in $\mathbb{R}$, i.e. $\sqrt{2} \in \mathbb{R}$
        \item (CCC) If a sequence of real numbers converges, its limit is in $\mathbb{R}$, 
this statement is equivalent to CA. 
    \end{enumerate}
\end{remark}

\begin{theorem}[Archimedean Property]
    $\mathbb{N}$ is unbounded from above in $\mathbb{R}$.
\end{theorem}

\begin{proof}
    Suppose $\mathbb{N}$ is bounded from above by $m \in \mathbb{R}$. 
    \begin{align*}
        \implies & \exists \sup \mathbb{N} \in \mathbb{R} \text{ by CA} \\
        \implies &m^{\prime} := \sup \mathbb{N} \text{ is LUB} \\
        \implies & m^{\prime} -1 \text{ is not an upper bound} \\
        \implies & \exists n \in \mathbb{N} \text{ s.t. } m^{\prime} - 1 < n \\
        \implies & m^{\prime} < n + 1 \in \mathbb{N}
    \end{align*}

    This contradicts with the assumption that $m^{\prime}$ is the supremum.
\end{proof}

\begin{proposition}[Equivalent statements to Archimedean]
    \hfill
    \begin{enumerate}
        \item $\forall x > 0, y \in \mathbb{R}$, $\exists n \in \mathbb{N}$ s.t. $nx > y$
        \item $\forall x > 0, \exists n \in \mathbb{N}$ s.t. $0 < \frac{1}{n} < x$
    \end{enumerate}
\end{proposition}

\begin{proof}
    \begin{enumerate}
    \item AP $\implies$ (1): Consider $\frac{y}{x} \in \mathbb{R}$. By AP, 
$\exists n \in \mathbb{N}$ s.t. $n > \frac{y}{x}$. Hence $nx > y$. 
    \item (1) $\implies$ (2): Take $y = 1$, then $\exists n \in \mathbb{N}$ s.t. 
$nx > 1 \implies x > \frac{1}{n} > 0$. 
    \item (2) $\implies$ (AP): By contrapositive. WTS $\mathbb{N}$ bounded from above 
$\implies \exists x \in \mathbb{R}, x > 0, \forall n \in \mathbb{N}, \frac{1}{n} > x$. 
\medskip 

Since $\mathbb{N}$ bounded, let $n_0$ denote an upper bound. Then let $x_0 = \frac{1}{n_0} > 0$. 
We have for arbitrary $n$
\begin{align*}
    n < n_0  \implies \frac{1}{n} > \frac{1}{n_0}  = x_0 
\end{align*}
    \end{enumerate}
\end{proof}
